\section{Discussion: memory systems need a verification scheduler}

The results describe a control problem that current memory architectures do not
expose. A retrieval layer answers \emph{which records are relevant to what the
agent is doing}. Nothing answers \emph{which records are worth spending scarce
provenance checks on}. In our environment those two rankings are close to
opposites: the belief most costly if wrong is the one outside the current plan,
and that is exactly the one the budget does not reach.

Making the second ranking explicit needs metadata the first does not: how many
consolidation generations a record has passed through, when its provenance was
last actually followed, how broad its scope claim is relative to the evidence
behind it, and which future contexts could make it load-bearing. Study 4 gives a
concrete signal for the third: a record that retains a number while losing its
scope or its prohibition has drifted in the direction that matters, and the drift
is detectable by comparing generations, which a consolidation pipeline is in a
position to do and an agent reading the final summary is not.

We are wary of the term \emph{verification debt}, but the accounting it suggests
is the useful part. A durable belief that survives many consolidations without
its provenance ever being followed accumulates risk that no single session sees,
because each session rationally spends its budget on its own plan. Study 1 shows
that behaviour is not a defect of any one model --- it is what every model we
tested does, and it is locally correct. The fix is therefore unlikely to be a
better-behaved agent. It is a scheduler outside the agent that can spend
verification against the store's accumulated exposure rather than against today's
decision.

\paragraph{Steering works, partially.} One system-prompt sentence naming no
memory and no direction moves allocation substantially (Table~\ref{tab:unified},
odds ratio \ORTriage), and Study 2a shows that this reaches behaviour one decision
later through what it caused the agent to recover. But it does not equalise
models, and in our replay the model with the weakest allocation still walks into
the corrupted direction a quarter of the time under the instruction. Prompt-level
triage is worth writing and is not a substitute for scheduling.

\section{Limitations}

\textbf{One scenario.} Every episode instantiates the same synthetic
growth-agent situation with six memories and five actions. We deliberately did
not add a second domain: built in the time available it would have been a
paraphrase of the first, and generality claimed from two structurally identical
scenarios is worse than generality not claimed. Domain replication and larger
memory pools are the obvious next step and we have not taken it.

\textbf{Small memory pool, perfect archive.} Six records, explicit ids, and a
provenance store that always answers correctly and instantly. Real stores are
larger, slower, sometimes missing, and sometimes themselves stale. We expect all
of these to sharpen the allocation problem, but we have not shown it.

\textbf{Study 1 is a re-analysis.} Its episodes were collected earlier under
their own pre-registration. The confounds we identified in its clean arm
(Study 3) and in its two-phase follow-up (Study 2) are why Studies 2--4 exist;
where Study 1's reported interpretation and our new evidence conflict, we have
said so in the section concerned rather than restating the earlier reading.

\textbf{Causal scope.} Studies 2a and 2b identify the effect of \emph{holding}
recovered provenance on the next decision. Neither identifies why models
allocate as they do, and nothing here tests whether allocation is stable across
tasks or sessions --- so we avoid calling it a trait.

\textbf{Intent and allocation are measured in the same response.} A model's
first-pass intended action and its verification request come from one structured
output. An agent constructing a coherent answer may name the memory behind the
action it is about to state, which would produce the separation we report without
any allocation policy behind it. The target-swap arm and Study 2 are what argue
against that reading --- the same corruption is caught when moved onto the
intended path, and what the budget recovered changes behaviour a decision later
--- but a design that elicited the allocation before the plan would settle it and
ours does not.

\textbf{Two exploratory analyses.} The intent-misaligned stratum in Study 3 was
chosen after seeing that the arms differed in intent (\S\ref{sec:qualifier}),
and the split of Study 4's corrupted memory into ``states the outcome'' and
``states a prohibition'' uses keyword lists written after reading the
generations. Neither was pre-registered. Both are reported as exploratory, and
the registered contrasts they sit beside are reported unchanged.

\textbf{The Study 4 scorer is one of the models under test.} Survival is judged
by a model from the same set that produced the consolidations. It passes the
false-positive control (\FiveFalsePosAll{} on records carrying no qualifier) and
identifies the unqualified sources correctly in every case, but an independent
judge would be better and we did not use one.

\textbf{Behaviour, not harm.} We measure which direction the agent chooses and
at what scale. Under a strict pre-registered harm rule requiring an unhedged,
unverified, corrupted-evidence rollout, no episode in this work or its
predecessor qualifies: models hedge verbally while still choosing the direction.
That is itself a warning about self-reported uncertainty as a safety metric, and
it means our outcome is exposure, not damage.

\section{Conclusion}

Provenance links do not make an inherited memory store safe; following them
does, and following is budgeted. We find that agents spend that budget on the
beliefs behind the plan they already have --- with no counterexample in
\ExpOneN{} episodes --- and that what the budget happens to recover causally
determines what they do when the situation later shifts, by roughly ninety
percentage points under two independent randomised designs. We also find that a
silently deleted caveat attracts less scrutiny than an irrelevant hedge, and
that real consolidation mostly does not delete the negative finding --- it strips
scope and normative force while keeping the numbers that make the lesson look
well-evidenced. Retrieval relevance and verification priority are different
rankings. Memory systems currently compute only the first.
